% Options for packages loaded elsewhere
\PassOptionsToPackage{unicode}{hyperref}
\PassOptionsToPackage{hyphens}{url}
%
\documentclass[
  10pt,
]{article}
\usepackage{amsmath,amssymb}
\usepackage{iftex}
\ifPDFTeX
  \usepackage[T1]{fontenc}
  \usepackage[utf8]{inputenc}
  \usepackage{textcomp} % provide euro and other symbols
\else % if luatex or xetex
  \usepackage{unicode-math} % this also loads fontspec
  \defaultfontfeatures{Scale=MatchLowercase}
  \defaultfontfeatures[\rmfamily]{Ligatures=TeX,Scale=1}
\fi
\usepackage{lmodern}
\ifPDFTeX\else
  % xetex/luatex font selection
  \setmainfont[]{DejaVu Serif}
  \setmonofont[]{DejaVu Sans Mono}
\fi
% Use upquote if available, for straight quotes in verbatim environments
\IfFileExists{upquote.sty}{\usepackage{upquote}}{}
\IfFileExists{microtype.sty}{% use microtype if available
  \usepackage[]{microtype}
  \UseMicrotypeSet[protrusion]{basicmath} % disable protrusion for tt fonts
}{}
\makeatletter
\@ifundefined{KOMAClassName}{% if non-KOMA class
  \IfFileExists{parskip.sty}{%
    \usepackage{parskip}
  }{% else
    \setlength{\parindent}{0pt}
    \setlength{\parskip}{6pt plus 2pt minus 1pt}}
}{% if KOMA class
  \KOMAoptions{parskip=half}}
\makeatother
\usepackage{xcolor}
\usepackage[margin=0.7in]{geometry}
\usepackage{longtable,booktabs,array}
\usepackage{calc} % for calculating minipage widths
% Correct order of tables after \paragraph or \subparagraph
\usepackage{etoolbox}
\makeatletter
\patchcmd\longtable{\par}{\if@noskipsec\mbox{}\fi\par}{}{}
\makeatother
% Allow footnotes in longtable head/foot
\IfFileExists{footnotehyper.sty}{\usepackage{footnotehyper}}{\usepackage{footnote}}
\makesavenoteenv{longtable}
\setlength{\emergencystretch}{3em} % prevent overfull lines
\providecommand{\tightlist}{%
  \setlength{\itemsep}{0pt}\setlength{\parskip}{0pt}}
\setcounter{secnumdepth}{-\maxdimen} % remove section numbering
\emergencystretch=3em
\sloppy
\ifLuaTeX
  \usepackage{selnolig}  % disable illegal ligatures
\fi
\usepackage{bookmark}
\IfFileExists{xurl.sty}{\usepackage{xurl}}{} % add URL line breaks if available
\urlstyle{same}
\hypersetup{
  hidelinks,
  pdfcreator={LaTeX via pandoc}}

\author{}
\date{}

\begin{document}

\section{Follow-up experiments for twist routes
1--3}\label{follow-up-experiments-for-twist-routes-13}

Date: 2026-10-22. \textbf{Design only: E18--E29 are proposed and have
not run.} These twelve proposals extend
\href{https://github.com/aburan28/crypto/pull/2380}{PR \#2380}. They
preserve every open requirement in the \href{../EXPERIMENTS.md}{original
E00--E17 program} and the \href{../reports/EXECUTION_RESULTS.md}{bounded
execution report}. No new computation, lease, funding transaction or
scheduled task is created.

The first scientific questions are actual P-521/toy 3-gluing (E18/E19),
coupled Frobenius-component restrictions (E21), and mixed-helper
elimination (E23). E25 and E27 close solver-reference and accounting
gaps needed to judge those routes. E24 and E29 close independent
arithmetic and concrete proof gaps. E20/E22/E26/E28 widen the search
without assuming a gain. All are retained; priority orders work and does
not cancel a proposal.

\subsection{What the current evidence leaves
open}\label{what-the-current-evidence-leaves-open}

The
\href{https://github.com/aburan28/crypto/tree/9dff98b15fd91887d41a73419a03abe17f9e24ca/research/twist_routes_1_3_20261021}{frozen
source snapshot} is the evidence base. Its module filter names P-521 as
a 3-torsion candidate; the quotient, polarization and maps are
unconstructed. Its cover arithmetic has native class checks but no
independent arbitrary-class replay. Its solver panel lacks an applicable
signature baseline and complete-basis certificates for every reported
partial output. Its IC comparisons leave material work unpriced. None
establishes a calibrated twist-assisted IC speedup or novelty.

The Frobenius separation argument continues to apply to adding pure
twist points to the original factor base. The collection hypotheses
below test the cost of finding verified projected rational relations;
they do not assume that an auxiliary representation creates a new group
relation. Geometry, construction, solver-stage improvement and complete
IC improvement are separate conclusions with separate certificates.

\subsection{Fixed design rules}\label{fixed-design-rules}

\begin{itemize}
\tightlist
\item
  \textbf{Inputs:} reuse the exact
  \href{../execution-v1/fixtures.json}{fixture manifest}. Its full
  SHA-256 is pinned in the \href{experiments.json}{machine matrix}. The
  ten ordinary fixtures are the two structural models and the eight
  ladder models. The three even/composite controls keep their original
  roles. Train on both structural models plus ladder replicate 0 at each
  size; hold out ladder replicate 1 at each size. This is a fixed
  retrospective split of already public curves, not a claim that those
  curves were unseen.
\item
  \textbf{New targets:} for experiments requiring new targets, labels
  0--15 are training, 16--31 validation, and 32--63 fresh final E28
  evaluation. A native generator must domain-separate with
  \texttt{twist-followup-20261022-v1}, the full curve ID, experiment ID
  and label, then rejection-sample a nonzero scalar using SHA-256. The
  exact serialization and generator revision must be frozen in the
  runnable protocol. Known scalars are verifier-only; collectors and
  selectors receive points. E25/E27 explicitly reuse their historical
  inputs.
\item
  \textbf{Controls:} same projected base, targets, sign conventions,
  arity, subgroup and target-span stopping rule. Compare the best
  verified ordinary IC and strongest applicable external solver, with a
  matched rho reference. Setup, failed attempts, projection,
  deduplication, verification, linear algebra and target recovery all
  count. Cold one-target cost is primary.
\item
  \textbf{Readiness:} native Rust implementations, exact
  inputs/argv/source pins, independent checkers, checkpointing and a
  committed hashed protocol are required before a card becomes ready.
  Every card is blocked and unfunded; its failing sentinel is not a
  launch command. Proposed caps below are limits, not runtime estimates
  or resource reservations.
\item
  \textbf{Decisions:} identical verified output with a cost reduction
  against the frozen best IC baseline is the improvement test. A stage
  result stays a stage result. Toy correctness, necessary filters, more
  lifts, lower degree, an energy statistic or ideal equality alone do
  not pass it. Preserve all failures and inconclusive cells. A bounded
  failure does not prove nonexistence. Do not use heldout outcomes to
  redesign the same evaluation.
\item
  \textbf{Formal claims:} this document proposes obligations; it asserts
  no new algebraic identity. Any executed round claiming one owes an
  identity certificate and the repository's Lean/source-bridge evidence.
\end{itemize}

\subsection{Priority and ownership}\label{priority-and-ownership}

P0 closes a comparison prerequisite; P1 tests a route/bridge; P2
broadens or diagnoses it. All statuses below are \textbf{proposed,
blocked, not started}.

\begin{longtable}[]{@{}
  >{\raggedright\arraybackslash}p{(\columnwidth - 8\tabcolsep) * \real{0.0700}}
  >{\raggedright\arraybackslash}p{(\columnwidth - 8\tabcolsep) * \real{0.5300}}
  >{\raggedright\arraybackslash}p{(\columnwidth - 8\tabcolsep) * \real{0.1000}}
  >{\raggedright\arraybackslash}p{(\columnwidth - 8\tabcolsep) * \real{0.0700}}
  >{\raggedright\arraybackslash}p{(\columnwidth - 8\tabcolsep) * \real{0.2300}}@{}}
\toprule\noalign{}
\begin{minipage}[b]{\linewidth}\raggedright
ID
\end{minipage} & \begin{minipage}[b]{\linewidth}\raggedright
Proposal
\end{minipage} & \begin{minipage}[b]{\linewidth}\raggedright
Priority
\end{minipage} & \begin{minipage}[b]{\linewidth}\raggedright
Route
\end{minipage} & \begin{minipage}[b]{\linewidth}\raggedright
Depends on
\end{minipage} \\
\midrule\noalign{}
\endhead
\bottomrule\noalign{}
\endlastfoot
\hyperref[e18]{E18} & Actual 3-torsion and pairing certificates & P1 & 2
& existing evidence \\
\hyperref[e19]{E19} & Indecomposable gluing, a genus-2 curve and
explicit maps & P1 & 2 & E18 \\
\hyperref[e20]{E20} & Other small gluing primes and actual torsion
modules & P2 & 2 & existing evidence \\
\hyperref[e21]{E21} & Restrictions coupling the two Frobenius components
& P1 & 1 & existing evidence \\
\hyperref[e22]{E22} & Projected additive energy and relation-yield
diagnostics & P2 & 1 & E21 \\
\hyperref[e23]{E23} & Mixed extension helpers and partial-relation
elimination & P1 & 1 & E21, E29 \\
\hyperref[e24]{E24} & Base-field Jacobian arithmetic and independent
arbitrary-class replay & P1 & 2 & existing evidence \\
\hyperref[e25]{E25} & Applicable homogeneous signature baseline with
complete output certificates & P0 & 3 & existing evidence \\
\hyperref[e26]{E26} & Frobenius blocks, hybrid elimination and bounded
invariant encodings & P2 & 3 & E25 \\
\hyperref[e27]{E27} & Complete primitive accounting and a cold
one-target replay & P0 & 3 & existing evidence \\
\hyperref[e28]{E28} & Heldout family selection and an honest solver
portfolio & P2 & 3 & E21, E22, E26, E27 \\
\hyperref[e29]{E29} & Concrete separation and projected target-span
certificates & P1 & 1 & existing evidence \\
\end{longtable}

The \href{experiments.json}{machine matrix} holds the same designs and
backlog IDs. The \href{PROGRAM.svg}{dependency diagram}, editable
\href{PROGRAM.dot}{Graphviz source} and \href{REPORT.pdf}{PDF report}
accompany this report. Every new proposal has its own card; the three
original route cards remain the parent trackers.

\subsection{E18}\label{e18}

\textbf{Actual 3-torsion and pairing certificates.} P1; route 2;
construction bridge.
\href{../../../backlog/XP-20261022-twist-actual-three-torsion/README.md}{Blocked
backlog card}.

\textbf{Question.} Do the four companion-matrix candidates lift to
explicit Galois-equivariant anti-isometries on actual 3-torsion,
including for public P-521 parameters?

\textbf{Why now.} The public module panel selected P-521 and three toy
curves, but constructed no torsion basis or Weil pairing. This is the
first missing certificate, before any quotient or performance question.

\textbf{Hypothesis, untested.} At least one matrix candidate can be
represented by actual torsion points with the claimed Frobenius action
and inverse Weil pairing. A positive answer establishes only this
bridge.

\textbf{Inputs and finite grid.} Use the three positive fixtures at
p=211, 751 and 10111 from PUBLIC\_3\_MODULE\_ANALYSIS.json, plus public
P-521. Keep all seven negative ordinary toy rows and the four negative
NIST rows as filter controls. Use the exact registered models, twist
parameter and public orders; no NIST target logarithm is supplied.

\textbf{Procedure.} Construct the 3-division polynomial and a declared
extension field. Search extension degrees 1 through 8 in increasing
order; certify irreducibility and all coordinate conversions. Give
independent order-3 points P,Q, their primitive Weil pairing, and both
Frobenius matrices in explicitly related symplectic bases. Replay every
candidate map and check its graph is Galois stable and reverses the
pairing. Native BigInt field/torsion support is a prerequisite for
P-521; the toy u128 adapter is insufficient.

\textbf{Controls.} A matrix with the wrong determinant, a deliberately
conjugated wrong Frobenius matrix, dependent torsion points and each
negative module row must fail the appropriate check. An empty matrix set
is not passed to a point constructor as a success.

\textbf{Required outputs.} Exact field polynomials and bases, all
torsion coordinates, basis-change matrices, pairing values, Frobenius
images and independent verification transcripts. Record construction
operations and peak memory, separately from the certificate conclusion.

\textbf{Decision and falsifier.} A verified anti-isometry enables E19. A
mismatch is a specific module-to-geometry defect to resolve.
Extension/resource exhaustion leaves that candidate unresolved; it does
not refute existence.

\textbf{Bounded pilot.} Proposed pilot: one CPU, 4 GiB per process,
1,800 seconds per positive curve, extension degree at most 8. Preserve
each attempted field and each failure. This is a cap, not an estimated
runtime.

\textbf{Dependencies and blockers.} Existing frozen evidence only.
Actual torsion/pairing implementation, native BigInt adapter and an
independent certificate checker are unimplemented.

\textbf{Evidence pointers.}
\href{../../../research/twist_routes_1_3_20261021/reports/execution-20261022T0032Z/PUBLIC_3_MODULE_ANALYSIS.json}{PUBLIC
3 MODULE ANALYSIS.json};
\href{../../../research/twist_routes_1_3_20261021/PUBLIC_3_MODULE_PROTOCOL_V1.md}{PUBLIC
3 MODULE PROTOCOL V1.md}.

\subsection{E19}\label{e19}

\textbf{Indecomposable gluing, a genus-2 curve and explicit maps.} P1;
route 2; geometric hypothesis.
\href{../../../backlog/XP-20261022-twist-three-gluing-geometry/README.md}{Blocked
backlog card}.

\textbf{Question.} Does an E18 anti-isometry yield an indecomposably
principally polarized quotient with an explicit same-field genus-2 model
and subgroup transfer?

\textbf{Why now.} A compatible Frobenius polynomial is not a curve. A
curve without maps is not an ECDLP transfer. The original positive
three-gluing requirement remains unfulfilled.

\textbf{Hypothesis, untested.} One of the certified toy anti-isometries
yields a usable genus-2 quotient. Public P-521 is a separate
construction-only extension after the toy implementation is verified.

\textbf{Inputs and finite grid.} Process every E18-certified
anti-isometry on the three positive toy models, up to the four
candidates reported per model. Retain every rejected candidate. Attempt
at most the four public P-521 candidates only after the same formulas
and checker pass the toy stage.

\textbf{Procedure.} Construct the graph quotient of E times its twist,
certify the induced principal polarization and geometric
indecomposability, descend a smooth genus-2 curve to the declared base
field, and provide quotient/pullback/norm maps with their actual
degrees. Prove each composition identity and the exact kernel
intersection with the source prime subgroup; normalize recovery by the
actual coprime degree. Do not infer degree or injectivity from the
numeral 3 alone.

\textbf{Controls.} Include a documented reducible/product-polarization
example, a non-anti-isometric graph and identity/ramification cases.
Compare independent point counts and divisor-map replay. A product
principally polarized surface, matching j-invariant or characteristic
polynomial cannot pass the smooth-curve certificate.

\textbf{Required outputs.} Polarization and kernel certificates, curve
equation, smoothness/genus/infinity checks, explicit rational maps,
exceptional charts, exact subgroup order, scalar recovery and all
failures. Initially measure construction and map costs only; no
Jacobian-IC improvement follows from existence.

\textbf{Decision and falsifier.} A fully checked toy model becomes an
E24 arithmetic input and a future destination-solver input. A failed
anti-isometry closes only that candidate. P-521 construction success
remains a public-parameter result, with no NIST DLP execution.

\textbf{Bounded pilot.} Proposed pilot: one CPU, 4 GiB and 3,600 seconds
per anti-isometry. Stop that cell on a singular model, failed descent,
failed map identity or resource cap; retain the failed certificate.

\textbf{Dependencies and blockers.} E18 Explicit
gluing/polarization/model/map constructors and independent replay are
missing.

\textbf{Evidence pointers.}
\href{../../../research/twist_routes_1_3_20261021/EXPERIMENTS.md}{EXPERIMENTS.md};
\href{../../../research/notes/index-calculus/RESEARCH_COVER_DECOMPOSITION_LEDGER.md}{RESEARCH
COVER DECOMPOSITION LEDGER.md}.

\subsection{E20}\label{e20}

\textbf{Other small gluing primes and actual torsion modules.} P2; route
2; applicability search.
\href{../../../backlog/XP-20261022-twist-higher-gluing-filters/README.md}{Blocked
backlog card}.

\textbf{Question.} Do gluing primes 5, 7, 11, 13 or 17 expose additional
public-parameter candidates that the 3-torsion filter misses?

\textbf{Why now.} The 3-torsion panel tests one gluing degree. A
negative row there does not exclude another prime, while a positive
matrix row still needs the geometric checks in E18/E19.

\textbf{Hypothesis, untested.} Some frozen curve/twist pairs have
suitable small-prime torsion modules at another degree. This proposes an
applicability census, not a new gluing construction.

\textbf{Inputs and finite grid.} All ten ordinary toy fixtures and all
five public NIST P-curves, crossed with primes 5,7,11,13,17. Use the
original public orders to reconstruct trace and characteristic. Reject a
cell whose gluing prime is the characteristic or divides the selected
source subgroup order.

\textbf{Procedure.} Derive and independently check the actual necessary
Frobenius conditions, then solve the anti-intertwiner linear system and
determinant condition. Separate companion-matrix results from actual
torsion action, especially scalar/repeated-eigenvalue cases. For each
prime, select the first two positive toy fixture IDs in lexical order,
if any, for explicit torsion/pairing replay; retain a no-positive status
rather than generating a favorable replacement fixture.

\textbf{Controls.} Wrong-trace and wrong-determinant matrices, the known
3-torsion panel as a regression control, characteristic-prime rejection
and a repeated-eigenvalue synthetic module control. Matrix counts are
not geometric certificates.

\textbf{Required outputs.} A complete 75-cell ordinary/public filter
table, all exclusions, module types, candidate matrices and bounded
actual-torsion replay. Actual quotients at the larger degrees remain a
separately frozen extension of E19.

\textbf{Decision and falsifier.} A verified positive adds a specifically
named construction candidate. An all-negative table excludes only this
finite prime set. Rising map/construction costs are reported before
expanding gluing degree.

\textbf{Bounded pilot.} Proposed pilot: one CPU, 4 GiB, 15 seconds per
matrix cell; at most two explicit toy torsion constructions per prime
with a 1,800-second cap each. No public large-prime torsion construction
is included in this pilot.

\textbf{Dependencies and blockers.} Existing frozen evidence only. A
generalized native module enumerator and actual-torsion checker are not
yet implemented.

\textbf{Evidence pointers.}
\href{../../../research/twist_routes_1_3_20261021/public-three-gluing.gp}{public-three-gluing.gp};
\href{../../../research/twist_routes_1_3_20261021/reports/execution-20261022T0032Z/PUBLIC_3_MODULE_ANALYSIS.json}{PUBLIC
3 MODULE ANALYSIS.json}.

\subsection{E21}\label{e21}

\textbf{Restrictions coupling the two Frobenius components.} P1; route
1; collection hypothesis.
\href{../../../backlog/XP-20261022-twist-coupled-component-bases/README.md}{Blocked
backlog card}.

\textbf{Question.} Can a predicate involving both the rational and
anti-rational components find useful relations more cheaply than direct
collection on exactly the same projected base?

\textbf{Why now.} Unrestricted extra twist fibers duplicate rational
traces. A predicate coupling the components is a concrete way to test a
different search representation while preserving the original separation
argument.

\textbf{Hypothesis, untested.} A coupled predicate makes decomposition
easier even when the final rational relation set is unchanged. Merely
increasing witness multiplicity does not support this hypothesis.

\textbf{Inputs and finite grid.} Use the ten ordinary fixtures and the
fixed training/holdout split above. For general extension points P
define \texttt{T=P+Frob(P)} and \texttt{U=P-Frob(P)}. Test
\texttt{x(T)+lambda*x(U)\ in\ W\_h} and \texttt{x(T)*x(U)\ in\ W\_h},
with lambda=1,2,3 and h=2,4,8; the product predicate is one family
independent of lambda, giving 12 distinct configurations. Here x(U) uses
the embedded twist coordinate and \texttt{W\_h=\{0,...,h-1\}} in Fp.

\textbf{Procedure.} First enumerate all extension points at
p=59,211,251. On larger fixtures use at most 2\^{}16 deterministic
public point proposals per predicate. Record every trace fiber and
construct signed projected bases of sizes 8,12,16 in canonical encoding
order. Compare the auxiliary system, direct rational collection on
exactly that base, a rational elimination of the auxiliary predicate,
and a matched random base. Handle T=O or U=O as separately enumerated
branches. Freeze at most two configurations using training data before
opening holdout outcomes.

\textbf{Controls.} Pure-twist helpers, uncoupled predicates with matched
acceptance rate, duplicated witnesses, shuffled auxiliary labels, exact
same projected base and the strongest ordinary decomposition method.
Unsupported or undersized bases remain failed cells.

\textbf{Required outputs.} Exact projected sets, fiber histograms,
canonical signed rows, target-row span, verified scalar recovery and
complete construction/collection/projection costs. Distinguish a new row
from a cheaper way to find an existing row.

\textbf{Decision and falsifier.} Only a lower verified cost per
independent projected row on heldout data justifies a solver/pipeline
extension. Equal projected rows with no cost decrease close the tested
predicate. No result bypasses Frobenius separation.

\textbf{Bounded pilot.} Proposed pilot: one CPU, 4 GiB; at most 2\^{}20
counted group operations per relation cell and 60 seconds as a separate
fail-safe. All construction and failed cells count. Exhaustive and
sampled populations are reported separately.

\textbf{Dependencies and blockers.} Existing frozen evidence only.
Coupled-family constructors, complete counters and independent
projected-row replay are required.

\textbf{Evidence pointers.}
\href{../../../research/twist_routes_1_3_20261021/execution-v1/fixtures.json}{fixtures.json};
\href{../../../research/twist_routes_1_3_20261021/reports/EXECUTION_RESULTS.md}{EXECUTION
RESULTS.md}.

\subsection{E22}\label{e22}

\textbf{Projected additive energy and relation-yield diagnostics.} P2;
route 1; structural diagnostic.
\href{../../../backlog/XP-20261022-twist-projected-energy/README.md}{Blocked
backlog card}.

\textbf{Question.} Do coupled trace images have an additive structure
that predicts more useful decompositions after cardinality and fiber
multiplicity are controlled?

\textbf{Why now.} Counting more lifts is inconclusive. Projected
additive energy and actual target decomposition histograms distinguish a
structured rational image from witness duplication.

\textbf{Hypothesis, untested.} A frozen projected family has higher
useful decomposition yield than matched random subsets of the same
rational group. This is a diagnostic and supplies no unknown-log
algorithm by itself.

\textbf{Inputs and finite grid.} The E21 exact projected sets, existing
ordinary bases and sizes 8,12,16. Use all targets on the three
exhaustive pilot curves and the fixed public target split on the other
fixtures. Include 32 public matched-random-base seeds per size.

\textbf{Procedure.} Compute pair-sum multiplicities, additive energy,
three-term target decompositions and independent signed-row yield using
actual group operations. Optional character/Fourier coordinates may use
known toy logs only inside a separate diagnostic checker; those logs
must never enter the collector, selector or reported algorithm cost.
Preserve raw histograms and compare identities among configurations
rather than treating repetitions as independent data.

\textbf{Controls.} Matched base cardinality, sign quotient, support,
acceptance rate and construction budget; fiber-resampled copies of one
base; shuffled labels; interval-like and random rational bases. Test
whether observed target yield follows the full histogram rather than its
mean alone.

\textbf{Required outputs.} Exact toy distributions and heldout
yield/rank tables, counted diagnostic work and uncertainty across
independent random bases. Record any mismatch between energy and usable
target-row span.

\textbf{Decision and falsifier.} A robust structural difference selects
a specifically frozen family for further work. An energy increase
without independent-row or cost improvement is not a collection win. A
known-log diagnostic is never presented as an executable DLP shortcut.

\textbf{Bounded pilot.} Proposed pilot: one CPU, 4 GiB, at most 2\^{}20
pair/group operations per base and 60 seconds per histogram cell; an
incomplete histogram stays incomplete.

\textbf{Dependencies and blockers.} E21 Projected sets must exist;
diagnostic/collector separation and independent row checking are
unimplemented.

\textbf{Evidence pointers.}
\href{../../../research/twist_routes_1_3_20261021/reports/EXECUTION_RESULTS.md}{EXECUTION
RESULTS.md};
\href{../../../research/notes/index-calculus/RESEARCH_INDEX_CALCULUS_FACTOR_BASE.md}{RESEARCH
INDEX CALCULUS FACTOR BASE.md}.

\subsection{E23}\label{e23}

\textbf{Mixed extension helpers and partial-relation elimination.} P1;
route 1; collection hypothesis.
\href{../../../backlog/XP-20261022-twist-partial-relation-elimination/README.md}{Blocked
backlog card}.

\textbf{Question.} Can general extension-point helpers reduce collection
cost after every helper is eliminated, compared with the identical
rational partial-relation graph?

\textbf{Why now.} Pure twist helpers cannot contribute a new odd-order
source component. General extension points have a rational component and
permit a partial-relation experiment, but their projected variables must
be accounted for honestly.

\textbf{Hypothesis, untested.} An auxiliary representation collects
partial relations cheaply enough to pay for graph storage, elimination
and projection. The possible gain is computational; no new mixed-group
identity is assumed.

\textbf{Inputs and finite grid.} The ten ordinary fixtures, E21 bases of
sizes 8,12,16, and helper caps 16,32,64. Use relation arities 2 and 3.
Admit an r-primary projector only after checking the exact
extension-group order and required coprimality; otherwise emit a typed
unsupported cell.

\textbf{Procedure.} Verify each proposed relation in the actual
extension group. Build unknown columns only for nonzero r-primary
projections of general helpers; pure twist helpers are zero in that
source component. Compare zero-helper, pure-twist, general-helper and
direct-rational-helper arms. Apply the identical graph/Schur elimination
policy, stopping on target-row span with recoverable scalar, and replay
every final row in E. Neither known scalar witnesses nor hidden helper
logs may guide collection.

\textbf{Controls.} A deliberately fake twist-log column must be
rejected. Compare the same projected helper identities, same graph, same
target and same stopping rule; include graph cycles yielding only zero
rows and a standard rational large-prime-style control.

\textbf{Required outputs.} Partial rows and their exact group witnesses,
projector certificates, helper identities, graph components, elimination
transforms, nonzero independent final rows, scalar recovery and all
collection/storage/elimination costs.

\textbf{Decision and falsifier.} A heldout complete-cost improvement
supports this collector. A larger partial graph with unchanged useful
span does not. Pure-twist-only output must agree with the E29 separation
certificate.

\textbf{Bounded pilot.} Proposed pilot: one CPU, 4 GiB; cap each cell at
2\^{}20 group operations, 2\^{}16 stored partial rows and 120 seconds.
Record whichever cap fires and the unresolved target span.

\textbf{Dependencies and blockers.} E21, E29 Certified projectors, exact
helper identities, replayable graph elimination and full counters are
missing.

\textbf{Evidence pointers.}
\href{../../../research/twist_routes_1_3_20261021/reports/EXECUTION_RESULTS.md}{EXECUTION
RESULTS.md};
\href{../../../research/notes/index-calculus/RESEARCH_TWIST_IC_ROUTES_1_3.md}{RESEARCH
TWIST IC ROUTES 1 3.md}.

\subsection{E24}\label{e24}

\textbf{Base-field Jacobian arithmetic and independent arbitrary-class
replay.} P1; route 2; implementation and validation bridge.
\href{../../../backlog/XP-20261022-twist-base-field-jacobian/README.md}{Blocked
backlog card}.

\textbf{Question.} How much of the present cover cost comes from the Fp6
chart, and do arbitrary-class maps replay in an independent same-field
representation?

\textbf{Why now.} The ten covers passed native class tests in an
odd-degree Fp6 branch chart. PARI checked orders, not an independent
arbitrary-class group law or norm. That distinction remains open.

\textbf{Hypothesis, untested.} A balanced two-infinity Fp implementation
avoids avoidable representation cost while preserving the exact same
classes and maps. Faster arithmetic alone would be an
implementation-stage result.

\textbf{Inputs and finite grid.} All ten frozen selected cover models
and their explicit chart conversions. Replay all 3,551 classes on the
smallest Jacobian, the original 256 public sampled classes on each other
model, and all original map/associativity controls. Add E19 models only
after their curve/map certificates pass.

\textbf{Procedure.} Implement the actual balanced divisor conventions
over Fp, including both infinity points and exceptional branches. Supply
a checked conversion to the existing Fp6 chart and certify base-field
descent. Use an independently implemented divisor law or an applicable
external CAS for arbitrary-class addition, pullback and norm. Compare
native operation counts on identical classes, not point counts as a
substitute for divisor replay. Include degree-3/4 irreducible support
factorization, rather than only split support.

\textbf{Controls.} Origin, ramification, both infinity points, non-split
supports, inverses, repeated support and the entire smallest class
population. A mismatched infinity convention, wrong field conversion or
non-descended class must fail.

\textbf{Required outputs.} Explicit model and divisor conventions,
chart/field conversion certificates, independent class-map replay,
factorization coverage, stage counters, peak memory and unsupported
branches. Keep order certificates separate from group-law certificates.

\textbf{Decision and falsifier.} A verified base-field stage improvement
enables a correctly priced cover pipeline comparison. Missing
independent arithmetic or exceptional charts leaves transfer validation
partial. No destination decomposition solver is supplied merely by
changing charts.

\textbf{Bounded pilot.} Proposed pilot: one CPU, 4 GiB, 3,600 seconds
per fixture with per-class checkpoints. Preserve partial exhaustive
progress and failures; do not replace the required whole smallest
population with samples.

\textbf{Dependencies and blockers.} Existing frozen evidence only.
Balanced Fp arithmetic, independent arbitrary-class checker and complete
non-split factorization support are missing.

\textbf{Evidence pointers.}
\href{../../../research/twist_routes_1_3_20261021/runner/src/jacobian.rs}{jacobian.rs};
\href{../../../research/twist_routes_1_3_20261021/reports/EXECUTION_RESULTS.md}{EXECUTION
RESULTS.md}.

\subsection{E25}\label{e25}

\textbf{Applicable homogeneous signature baseline with complete output
certificates.} P0; route 3; solver-control closure.
\href{../../../backlog/XP-20261022-twist-homogeneous-solver-control/README.md}{Blocked
backlog card}.

\textbf{Question.} Can a same-characteristic signature solver be
compared fairly after homogenization, saturation and dehomogenization
are fully verified and priced?

\textbf{Why now.} The installed signature-mode attempt rejected
nonhomogeneous inputs and a diagnostic reported a different working
characteristic. Its output is not an applicable same-field baseline.

\textbf{Hypothesis, untested.} A supported homogeneous path or a
different applicable signature implementation can supply the missing
baseline. This is a solver-coverage hypothesis, not a predicted native
win.

\textbf{Inputs and finite grid.} Start with p=59,211,251, target seed 0
from the original solver fixtures, arities 2 and 3, all four original
encodings and grevlex/lex: 48 input systems. After passing
applicability/correctness, extend to the original 1,930 configured V3
cells with their original identities, not newly favorable replacements.

\textbf{Procedure.} First verify the mode and executable really support
the requested finite characteristic; if they do not, mark that backend
inapplicable and select a supported implementation. Construct the
homogenized ideal, saturate by the homogenizing variable where required,
and separately record components at infinity. Check dehomogenized ideal
membership both ways and replay all projected solutions. Certify a
claimed complete basis by every required S-pair reduction and generator
inclusion, or an independent complete-basis comparison. Ideal equality
alone does not certify a Gröbner basis.

\textbf{Controls.} Wrong-characteristic output, intentionally truncated
bases, lost affine branches, spurious infinity components, same-field
direct F4/q0 and a homogeneous positive control. Include all
encoding/decoding, saturation and solution-extraction work.

\textbf{Required outputs.} Exact transformed inputs, actual
field/mode/version, complete-basis or explicit incomplete status,
projected solution sets, operations, memory and all failures. Compare
with the fastest available applicable external implementation; name
unavailable proprietary references.

\textbf{Decision and falsifier.} Only verified equal tasks admit a stage
comparison. If no applicable signature implementation is available, that
baseline remains missing rather than being replaced by a
wrong-characteristic run.

\textbf{Bounded pilot.} Proposed pilot: one CPU, 4 GiB, 30 seconds per
engine/system with the same hard envelope for each compared engine.
Freeze timeout semantics and record extraction/checking separately
before extending the grid.

\textbf{Dependencies and blockers.} Existing frozen evidence only.
Applicable signature backend, homogenization/saturation adapter and
complete-basis checker are not yet pinned.

\textbf{Evidence pointers.}
\href{../../../research/twist_routes_1_3_20261021/SOLVER_PROTOCOL_V3.md}{SOLVER
PROTOCOL V3.md};
\href{../../../research/twist_routes_1_3_20261021/reports/EXECUTION_RESULTS.md}{EXECUTION
RESULTS.md}.

\subsection{E26}\label{e26}

\textbf{Frobenius blocks, hybrid elimination and bounded invariant
encodings.} P2; route 3; solver hypothesis.
\href{../../../backlog/XP-20261022-twist-block-elimination/README.md}{Blocked
backlog card}.

\textbf{Question.} Can splitting rational/anti-rational variable blocks
or eliminating a small helper block reduce certified solver work on the
same projected relation task?

\textbf{Why now.} Equivalent product/chain/Boolean/symmetric encodings
were checked only on the first solver grid. Structure-aware elimination
is a distinct algorithmic possibility, even when no new geometric
relation exists.

\textbf{Hypothesis, untested.} A bounded block elimination or hybrid
enumeration removes expensive intermediate terms more cheaply than
direct F4 on an equivalent task.

\textbf{Inputs and finite grid.} The ten original ordinary fixtures,
base sizes 8,12,16 and arities 2,3,4. Freeze four candidates: direct
addition chain; Frobenius block order; enumeration of one
at-most-eight-value helper followed by elimination; and one explicitly
derived first-pair invariant encoding. Use training targets first and
retain every failed encoding.

\textbf{Procedure.} Prove and exhaustively test projection equivalence
on the smallest fixtures, including signs, denominators and identity
branches. If field equations are added, declare the finite-field task
change and compare the same constrained task in every arm. Record
Macaulay dimensions/nonzeros, reductions, complete-basis certificates
and output lifting. Freeze variable orders and at most one selected
candidate before heldout targets; use E25 for an applicable external
reference.

\textbf{Controls.} The best direct rational low-degree chain, all
original encodings, identical finite bases and target rows,
wrong-sign/denominator negative systems, enumeration overhead and an A/A
run. Lower displayed degree or ideal equality without completeness
cannot pass.

\textbf{Required outputs.} Native systems and exact projection
certificates, complete/incomplete solver statuses, stage work, root
extraction, memory and cost per independent verified projected row.

\textbf{Decision and falsifier.} A verified heldout stage gain can
proceed to a complete pipeline comparison under E27. Otherwise record a
loss, tie or incomplete cell for each frozen encoding; no coordinate
reformulation is automatically called an advance.

\textbf{Bounded pilot.} Proposed pilot: one CPU, 4 GiB, 30 seconds per
solver cell; cap helper enumeration at eight values and at most four
encodings. All branch solves are charged to the same candidate cell.

\textbf{Dependencies and blockers.} E25 Candidate derivations, native
adapters, full completeness/lifting checks and counters are missing.

\textbf{Evidence pointers.}
\href{../../../research/twist_routes_1_3_20261021/runner/src/solver.rs}{solver.rs};
\href{../../../research/twist_routes_1_3_20261021/reports/EXECUTION_RESULTS.md}{EXECUTION
RESULTS.md}.

\subsection{E27}\label{e27}

\textbf{Complete primitive accounting and a cold one-target replay.} P0;
route 3; measurement closure.
\href{../../../backlog/XP-20261022-twist-complete-cost-accounting/README.md}{Blocked
backlog card}.

\textbf{Question.} Do any apparent IC stage or pipeline savings survive
complete symmetric pricing of construction, solver, projection,
verification and linear algebra?

\textbf{Why now.} The initial panel has unpriced high-level arithmetic,
symbolic work and extraction. Its ratios are sensitivity calculations
and cannot select a calibrated winning IC baseline.

\textbf{Hypothesis, untested.} No direction of improvement is assumed.
The experiment closes an accounting gap so that a win, loss or
unresolved comparison has a defensible unit.

\textbf{Inputs and finite grid.} First instrument the original 24 arms
on p=59,211,251. Then replay the full original ten-curve/eight-target
grid with all 1,920 attempts, including empty-base cells. Freeze the
best verified ordinary-IC baseline for each workload before evaluating
any new E21/E23/E26 pipeline.

\textbf{Procedure.} Instrument primitive field operations through all
reference and candidate code, including inversions with a frozen
algorithm, extension conversion, rejected candidates, F4/symbolic work,
roots, signs/deduplication, matrix work and target recovery. Derive any
GAE conversion from the declared counted reference law and expose its
assumptions. Keep raw counts, setup and online intervals. Apply the
0/0.1/1 GAE sensitivity band symmetrically to any remaining unpriced
units; if material units remain, the complete-cost conclusion remains
unresolved. Obtain admitted isolation and A/A stability before making
wall-time claims.

\textbf{Controls.} Identical-base ordinary IC, the strongest matched rho
available for each curve, a same-binary A/A arm, cold construction per
target and independently recomputed counters. An error is an attempted
cost, not a zero-cost solved workload.

\textbf{Required outputs.} Every phase and failed attempt, target-span
stopping and verified scalar, peak RSS, count definitions, complete or
explicitly incomplete costs, independent replay/aggregation and matched
rho reference. Preserve the original campaign as a separate frozen
version.

\textbf{Decision and falsifier.} A lower verified cost than the frozen
best IC baseline is an IC improvement within this workload. Beating rho
is a separate test. An unexplained sensitivity reversal or missing
material cost blocks a complete-cost claim.

\textbf{Bounded pilot.} Proposed pilot: one CPU and 4 GiB per arm, 120
seconds per one-target cell, checkpoint after each arm/target. Full-grid
runtime remains unestimated; pilot measurements are required before
scheduling it.

\textbf{Dependencies and blockers.} Existing frozen evidence only.
Complete counters, admitted calibration/isolation, independent
aggregation and a runnable frozen protocol are missing.

\textbf{Evidence pointers.}
\href{../../../research/twist_routes_1_3_20261021/E15_SPEC_V1.json}{E15
SPEC V1.json};
\href{../../../research/twist_routes_1_3_20261021/PIPELINE_PROTOCOL_V1.md}{PIPELINE
PROTOCOL V1.md};
\href{../../../research/twist_routes_1_3_20261021/reports/EXECUTION_RESULTS.md}{EXECUTION
RESULTS.md}.

\subsection{E28}\label{e28}

\textbf{Heldout family selection and an honest solver portfolio.} P2;
route 3; adaptive-search hypothesis.
\href{../../../backlog/XP-20261022-twist-heldout-selection/README.md}{Blocked
backlog card}.

\textbf{Question.} Can a selector using only public instance features
improve cold one-target IC after all search and selection costs are
included?

\textbf{Why now.} A large candidate grid can produce a favorable cell by
chance. A useful adaptive method must choose without seeing the target
logarithm or the heldout outcome.

\textbf{Hypothesis, untested.} A small frozen selector chooses a better
family/solver than one fixed global configuration often enough to pay
its own cost.

\textbf{Inputs and finite grid.} Only completed training outputs from
E21/E22/E26 and the complete counters from E27. Freeze at most four
candidate configurations and one selector. Final evaluation uses fresh
target labels 32 through 63 on the four ladder-replicate-1 holdout
curves; earlier labels 16 through 31 are validation, not fresh final
evidence.

\textbf{Procedure.} Use public coefficients, field size, base size,
degree/sparsity and budgeted pilot statistics as features. Freeze a
deterministic rule on the training curves and labels 0 through 15; hash
it before final targets are generated. Charge feature extraction, pilot
failures and any per-target tuning. Report cold one-target cost first;
any amortized reuse is a separate workload with an explicit target count
and setup interval.

\textbf{Controls.} Best frozen fixed IC configuration, matched rho, a
random selector with the same pilot budget, an A/A selector and a
hindsight oracle clearly labeled non-executable. Neither verifier-only
known logs nor final-target outcomes may enter the selector.

\textbf{Required outputs.} Selector source/hash and allowed features,
every candidate/pilot attempt, choices, failures, verified final scalars
and paired complete costs. Evaluate all final cells, including
unsupported choices; do not replace them after inspection.

\textbf{Decision and falsifier.} A complete-cost reduction on the sealed
final population supports an adaptive engineering improvement and
justifies independent replication of this frozen selector. A
training-only gain, oracle result or failed heldout transfer does not
support that conclusion.

\textbf{Bounded pilot.} Proposed pilot: one CPU, 4 GiB, at most four
configurations and two budgeted pilot probes per target; 120 seconds
total per selector cell, including probes. Final resource estimates
await E27.

\textbf{Dependencies and blockers.} E21, E22, E26, E27 Verified
candidates, complete counters, a frozen selector and unopened final
target set are required.

\textbf{Evidence pointers.}
\href{../../../research/twist_routes_1_3_20261021/execution-v1/fixtures.json}{fixtures.json};
\href{../../../research/twist_routes_1_3_20261021/reports/EXECUTION_RESULTS.md}{EXECUTION
RESULTS.md}.

\subsection{E29}\label{e29}

\textbf{Concrete separation and projected target-span certificates.} P1;
route 1; formal and structural bridge.
\href{../../../backlog/XP-20261022-twist-projection-span-certificate/README.md}{Blocked
backlog card}.

\textbf{Question.} Can the existing abstract Frobenius-separation
argument be connected to the actual finite-field embedding and to the
exact row transformations used by an auxiliary collector?

\textbf{Why now.} Nine Lean theorems cover abstract/conditional
statements, but the concrete field/twist/Frobenius and row-span bridges
remain open. These checks prevent a partial-relation matrix from
inventing source information.

\textbf{Hypothesis, untested.} For the stated odd-order and projector
hypotheses, pure-twist helper columns add no source target span; general
helper projections must be represented explicitly. This is a proposed
formalization of the existing obstruction, not a universal computational
lower bound.

\textbf{Inputs and finite grid.} The existing Lean project and exact
p=59 finite-field/twist embedding, then the three frozen even/composite
controls. Consume E23 relation matrices only after a separate group
replay; synthetic matrices with invalid twist-log columns are negative
controls.

\textbf{Procedure.} Prove the concrete twist embedding and Frobenius
identities, instantiate the odd-subgroup/projector hypotheses, and
connect each recorded matrix elimination to the projected rational row
module. Ship the relevant identity certificates and Lean proofs without
sorry or custom axioms. Search bounded even/torsion controls for
counterexamples to each deliberately weakened hypothesis. State
precisely which equalities are general proofs, exhaustive finite checks
or samples.

\textbf{Controls.} Remove the odd-order assumption, use a wrong
Frobenius sign, omit the factor-two normalization, inject a fake
twist-log unknown and supply a noninvertible projector denominator. Each
must fail verification or trigger the correctly typed exception.

\textbf{Required outputs.} Lean source and axiom audit, concrete
field/embedding certificates, exact matrix transforms, target-span
replay and negative controls. The proof must not conclude that
equivalent encodings have equal computational cost.

\textbf{Decision and falsifier.} A checked bridge constrains E23 and
closes part of the original proof map. Failure identifies the exact
missing hypothesis or implementation defect. It does not retire the
coupled-representation or geometric routes.

\textbf{Bounded pilot.} Proposed pilot: one CPU, 4 GiB, 600 seconds per
Lean module/build shard and exhaustive finite checks only on the
smallest frozen fixtures. Compilation failure or cap exhaustion leaves
the bridge open.

\textbf{Dependencies and blockers.} Existing frozen evidence only.
Concrete finite-field/Frobenius/twist and matrix-row-module formal
bridges are unimplemented.

\textbf{Evidence pointers.}
\href{../../../research/twist_routes_1_3_20261021/lean/README.md}{README.md};
\href{../../../research/twist_routes_1_3_20261021/execution-v1/fixtures.json}{fixtures.json};
\href{../../../docs/identity-certificates/README.md}{README.md}.

\subsection{Literature and novelty
gate}\label{literature-and-novelty-gate}

Summation polynomials, symmetry reductions, gluing, large-prime partial
relations, homogenization and adaptive solver portfolios are established
techniques. The possible contribution here is a verified construction or
measured method for these exact twist-assisted tasks, if one is found.
\textbf{Novelty has not been established.} These are hypotheses and
replication/ measurement gaps, not twelve new algorithms.

Before promoting a positive result, audit its exact construction,
assumptions, algorithm and complexity against at least these existing
source leads and the repository's
\href{../../notes/index-calculus/RESEARCH_COVER_DECOMPOSITION_LEDGER.md}{cover/decomposition
ledger}:

\begin{itemize}
\tightlist
\item
  \href{https://eprint.iacr.org/2004/031}{Semaev, summation polynomials,
  ePrint 2004/031}.
\item
  \href{https://www.iacr.org/archive/eurocrypt2014/84410158/84410158.pdf}{Faugere--Huot--Joux--Renault--Vitse,
  symmetrized summation polynomials, EUROCRYPT 2014}.
\item
  \href{https://eprint.iacr.org/2011/020}{Joux--Vitse, cover and
  decomposition, ePrint 2011/020}, with its extension-field assumptions
  kept explicit.
\item
  \href{https://arxiv.org/abs/1403.6911}{Broker--Howe--Lauter--Stevenhagen,
  genus-2 constructions}.
\item
  \href{https://nvlpubs.nist.gov/nistpubs/SpecialPublications/NIST.SP.800-186.pdf}{NIST
  SP 800-186, public parameters}.
\end{itemize}

These links are inherited source leads, not a fresh exhaustive
literature search. A future positive result must identify its closest
prior method and say whether it is a construction, implementation,
encoding or accounting improvement. Nothing here is a NIST DLP or
security claim.

\subsection{Delivery and unchanged
records}\label{delivery-and-unchanged-records}

The request to keep proposing experiments is addressed by twelve
concrete designs, a machine matrix, twelve blocked cards, this report,
its editable dependency diagram and PDF. No scientific experiment ran in
this follow-up. Validation is document/link/hash checking, native
backlog validation and the repository's required checks; those checks
are not experiment outcomes.

The prior protocols, raw sessions, receipts, result reports and PDFs
stay unchanged. The canonical catalog/browser and cover assessment
receive no new models. The checked canonical
views---\texttt{docs/index-calculus-scoreboard.html},
\texttt{docs/ic-leaderboard.html},
\texttt{docs/ic/progress-timeline.json}, the cover assessment diagrams
and the original study \texttt{PROGRAM.svg}---stay unchanged: this round
adds proposed dependencies, no verified curve, edge or measurement. No
new benchmark entry is warranted because no performance was measured.

\subsection{Proposed dependency
diagram}\label{proposed-dependency-diagram}

The following final PDF page contains the vector diagram. Every dashed
edge is a proposed dependency, not a verified transfer or measured
improvement. Construction and collector branches may be developed while
accounting and solver prerequisites are closed. E19 models can enter E24
only after their certificates pass; E20 positives require a new
degree-specific protocol.

\end{document}
